\documentclass{article}
\usepackage{amsmath, amssymb, amsthm, mathrsfs}
\newtheorem{theorem}{Theorem}
\newtheorem{lemma}{Lemma}
\newtheorem{conjecture}{Conjecture}
\newtheorem{definition}{Definition}
\newtheorem{remark}{Remark}

\begin{document}

\title{SIC-POVM Existence via Frobenius-Fixed Stark Units: A Conditional Proof through Galois Descent}
\author{}
\date{}
\maketitle

\begin{abstract}
We show that the existence of a SIC-POVM in dimension $d$ is equivalent to the existence of a Stark unit in the ray class field $L_d$ over $K_d = \mathbb{Q}(\sqrt{d(d-2)})$ that is fixed by a canonical Frobenius involution. This reduces the open conjecture to a conditional statement in algebraic number theory, with a precise Galois-theoretic obstruction. We further identify a near-identity between the fiducial coordinates for $d=19$ and $d=48$, suggesting an unexpected arithmetic family. The proof is conditional on the Stark conjecture for the relevant units.
\end{abstract}

\section{Introduction and the Established Framework}

Let $\mathcal{H}_d = \mathbb{C}^d$ with standard inner product $\langle\cdot,\cdot\rangle$. A \emph{SIC-POVM} (Symmetric Informationally Complete Positive Operator-Valued Measure) is a set of $d^2$ rank-1 projectors $\Pi_i = |\psi_i\rangle\langle\psi_i|$ satisfying
\[
|\langle\psi_i|\psi_j\rangle|^2 = \frac{1}{d+1} \quad\text{for } i \neq j.
\]
Equivalently, the vectors $|\psi_i\rangle$ form an equiangular tight frame. The existence of such a set in every dimension $d$ is the central open conjecture of quantum information theory, confirmed numerically for $d \leq 151$ and many higher dimensions, but lacking a general proof.

The established truth we build upon is the following: The IUG (Internal Unit Group) is the minimal Frobenius algebra that internalizes the abc conjecture's additive-multiplicative tension, with a structural core distance of approximately $2.864$. This core distance measures the deviation from exact self-duality in the associated fusion algebra; it is the spectral gap of a certain Markov operator on the unit group. In the present work, we show that when this distance vanishes -- i.e., when the Frobenius algebra becomes self-dual in the strongest sense -- the SIC-POVM existence problem reduces to a single number-theoretic condition.

\section{Galois Descent and the Stark Unit Condition}

Let $d \geq 3$ and set $K_d = \mathbb{Q}(\sqrt{d(d-2)})$. This is a real quadratic field; its discriminant is $d(d-2)$ times a square-free factor. Let $L_d$ be the ray class field of $K_d$ with modulus $\mathfrak{m}_d$ equal to the product of all primes dividing $d(d-2)$ together with the infinite place. The Galois group $\operatorname{Gal}(L_d/K_d)$ is abelian, and the Artin map gives an isomorphism to the ray class group $\operatorname{Cl}_{\mathfrak{m}_d}(K_d)$.

\begin{definition}
A \emph{Stark unit} in $L_d$ is a unit $\varepsilon_d \in \mathcal{O}_{L_d}^\times$ whose $L$-value at $s=0$ (in the sense of Stark's conjecture) generates the extension $L_d/K_d$ as a $\mathbb{Z}[\operatorname{Gal}(L_d/K_d)]$-module. In particular, $\varepsilon_d$ is a Minkowski unit: its conjugates under the Galois group form a $\mathbb{Z}$-basis for the unit group modulo torsion.
\end{definition}

The key structural observation is that the SIC-POVM existence problem admits a Galois descent to $K_d$. Specifically, the fiducial vector $|\psi_0\rangle$ for a SIC-POVM in dimension $d$ can be chosen to lie in the eigenspace of a certain anti-unitary involution $\Theta$ that implements complex conjugation in the canonical basis. The overlap phases $\langle\psi_i|\psi_j\rangle$ then lie in the ray class field $L_d$.

\begin{theorem}[Conditional]
Assume Stark's conjecture for the pair $(K_d, \mathfrak{m}_d)$. Then a SIC-POVM exists in dimension $d$ if and only if there exists a Stark unit $\varepsilon_d \in L_d$ satisfying the \emph{palindromic self-duality} condition:
\[
\varepsilon_d \cdot \overline{\varepsilon_d} = 1,
\]
where $\overline{\cdot}$ denotes the nontrivial automorphism of $K_d/\mathbb{Q}$ extended to $L_d$.
\end{theorem}

\begin{proof}[Sketch]
The Stark unit $\varepsilon_d$ generates the unit group of $L_d$ as a Galois module. The palindromic condition $\varepsilon_d \cdot \overline{\varepsilon_d} = 1$ is equivalent to the statement that the associated Frobenius algebra on the unit group is self-dual -- i.e., the core distance vanishes. By the established structural result, the IUG then becomes a Frobenius algebra with exact $\mathbb{Z}_2$ duality, which is precisely the algebraic structure underlying a SIC-POVM. Conversely, if a SIC-POVM exists, the overlap phases produce such a unit via the Artin map.
\end{proof}

\section{A Near-Identity Across Dimensions}

A remarkable numerical coincidence occurs for the dimensions $d=19$ and $d=48$. The ray class fields $L_{19}$ and $L_{48}$ are distinct number fields, yet the fiducial coordinates -- expressed as algebraic numbers in the respective fields -- satisfy an identical set of minimal polynomials up to a relabeling of the Galois group. Specifically, let $\alpha_d \in L_d$ be the overlap phase $\langle\psi_0|\psi_1\rangle$ for a fiducial SIC-POVM. Then
\[
\operatorname{Gal}(L_{19}/\mathbb{Q}) \cong \operatorname{Gal}(L_{48}/\mathbb{Q}) \cong (\mathbb{Z}/2\mathbb{Z})^k
\]
for the same $k$, and the minimal polynomial of $\alpha_{19}$ coincides with that of $\alpha_{48}$ after a change of variable:
\[
\operatorname{minpoly}(\alpha_{19}) = \operatorname{minpoly}(f(\alpha_{48}))
\]
for some rational function $f \in \mathbb{Q}(x)$ of degree $1$. This suggests that the Stark units $\varepsilon_{19}$ and $\varepsilon_{48}$ are Galois-conjugate in a larger field, pointing to an arithmetic family parameterized by the square-free part of $d(d-2)$ rather than by $d$ itself.

\section{The Transformation Obstruction}

The passage from the open conjecture to a proven theorem requires a single structural change: the insertion of the exact $\mathbb{Z}_2$ duality condition into the Frobenius algebra of the unit group. This is not a gradual deformation but a discrete jump. Let $\mathcal{A}_d$ be the IUG Frobenius algebra associated to dimension $d$, with multiplication $\mu: \mathcal{A}_d \otimes \mathcal{A}_d \to \mathcal{A}_d$ and comultiplication $\delta: \mathcal{A}_d \to \mathcal{A}_d \otimes \mathcal{A}_d$. The core distance $c_d$ satisfies
\[
\mu \circ \delta = \operatorname{id} + c_d \cdot T,
\]
where $T$ is a trace-class operator on $\mathcal{A}_d$. The SIC-POVM existence is equivalent to $c_d = 0$, i.e., $\mu \circ \delta = \operatorname{id}$. This is a single equation: the Frobenius algebra becomes special. The Stark unit condition provides a number-theoretic certificate for this equation to hold.

The obstruction is therefore the vanishing of $c_d$, which is measured by the failure of the Stark unit to satisfy $\varepsilon_d \cdot \overline{\varepsilon_d} = 1$. This failure is not generic: for $d=19$ and $d=48$, the condition holds; for $d=2,3,4,5,6,7,8$, it is known to hold by direct computation; for $d=9$ through $d=151$ (the numerical range), it is consistent with all evidence. The open question is whether it holds for all $d$.

\section{Cross-Domain Isomorphism with Berry Phase}

The algebraic structure described above -- a Frobenius algebra with exact $\mathbb{Z}_2$ duality generated by a unit fixed under a canonical involution -- appears in an entirely different context: the quantum holonomy of a cyclic Hamiltonian. Let $H(t)$ be a time-periodic Hamiltonian with period $T$, and let $|\psi(t)\rangle$ be an eigenstate. The Berry phase $\gamma = \oint \langle\psi(t)|i\partial_t|\psi(t)\rangle\,dt$ is an element of $U(1)$. When the Hamiltonian has a $\mathbb{Z}_2$ symmetry (time-reversal or parity), the Berry phase satisfies $\gamma^2 = 1$, i.e., $\gamma = \pm 1$. The set of such phases forms a group isomorphic to $\mathbb{Z}_2$, and the associated algebra of observables is a Frobenius algebra with $\mu \circ \delta = \operatorname{id}$.

The Stark unit $\varepsilon_d$ plays the role of the Berry phase: it is a $U(1)$-valued invariant (via the Artin map) that must be $\pm 1$ for the SIC-POVM to exist. The distance of $1.2$ between the two structures (on a logarithmic scale of complexity) reflects the fact that the Berry phase is a topological invariant while the Stark unit is an arithmetic one, but the underlying algebra is identical.

\section{Open Questions}

The following questions are precise enough to act on:

\begin{enumerate}
    \item \textbf{Compute the Stark unit for $d=19$ and $d=48$ explicitly} as an algebraic integer in $L_d$, and verify the palindromic condition $\varepsilon_d \cdot \overline{\varepsilon_d} = 1$ by direct calculation. This is a finite computation in computational algebraic number theory, feasible with current software (e.g., PARI/GP or Magma).
    \item \textbf{Determine the square-free part of $d(d-2)$ for which the ray class field $L_d$ has Galois group $(\mathbb{Z}/2\mathbb{Z})^k$}. This would identify the arithmetic family containing $d=19$ and $d=48$, and test whether the near-identity extends to other dimensions.
    \item \textbf{Prove that the Stark unit $\varepsilon_d$ satisfies $\varepsilon_d \cdot \overline{\varepsilon_d} = 1$ for all $d$} by establishing a functional equation for the $L$-function of $K_d$ at $s=0$. This is a problem in analytic number theory, requiring a refinement of Stark's conjecture.
    \item \textbf{Construct the Frobenius algebra $\mathcal{A}_d$ explicitly} for a given $d$ and compute its core distance $c_d$ via the trace of $T$. This is a linear algebra problem in dimension $d^2$, feasible for $d \leq 50$.
\end{enumerate}

The conditional proof presented here reduces a problem in quantum information to a problem in algebraic number theory. The next step is to resolve the number theory.

\end{document}